\section{The invariant double-cover classification}

The preceding sections studied several regular quotient presentations
inside one total graph. We now change viewpoint and classify
\(S_5\)-invariant binary double covers over the fixed base

\[
Y=L(P).
\]

These are different classification problems. A cohomology class records a
cover over a selected base, whereas a quotient frame records one covering
presentation inside a selected total graph.

\subsection{The fixed cohomology space}

The connected graph \(Y\) has fifteen vertices and thirty edges. Hence

\[
\dim_{\Ftwo} H^1(Y,\Ftwo)
=
|E(Y)|-|V(Y)|+1
=
16.
\]

The natural \(S_5\)-action on \(Y\) induces an action on this cohomology
space.

% STATUS: Certified computational input
% SOURCE: Project 41 packet 028m30i
% Exact repository path to be normalized in Section 9.

\begin{theorem}
\label{thm:fixed-cohomology}

The fixed subspace has dimension two:

\[
H^1(L(P),\Ftwo)^{S_5}
\cong
\Ftwo^2.
\]

Consequently there are exactly four \(S_5\)-invariant switching classes of
binary voltage assignments on \(L(P)\).
\end{theorem}

\subsection{Triangle and pentagon coordinates}

The cycle space of \(Y\) has dimension sixteen. The triangles of \(Y\)
span a ten-dimensional subspace. A certified pentagon orbit supplies the
remaining six dimensions, so the triangle and pentagon data together span
the full cycle space.

% STATUS: Certified computational input
% SOURCE: Project 41 packet 028m30i

On the fixed cohomology subspace, choose the two coordinate functionals

\[
t:
H^1(Y,\Ftwo)^{S_5}
\longrightarrow
\Ftwo
\]

and

\[
p:
H^1(Y,\Ftwo)^{S_5}
\longrightarrow
\Ftwo,
\]

given by evaluation on the certified triangle and pentagon test cycles.
The resulting map

\[
\kappa:
H^1(Y,\Ftwo)^{S_5}
\longrightarrow
\Ftwo^2,
\qquad
[\phi]\longmapsto(t([\phi]),p([\phi]))
\]

is an isomorphism.

For

\[
(a,b)\in\Ftwo^2,
\]

write

\[
\mathcal C_{ab}
\]

for the invariant switching class having coordinate \((a,b)\).

\subsection{The four invariant classes}

% STATUS: Certified computational input
% SOURCE: Project 41 cohomology-square and cover-order packets

The four classes have the certified properties shown in
Table~\ref{tab:invariant-cover-square}.

\begin{table}[h]
\centering
\begin{tabular}{lllll}
\toprule
Class &
Connected &
Triangles &
Bipartite &
Automorphism order\\
\midrule
\(\mathcal C_{00}\) & No  & Yes & No  & \(28800\)\\
\(\mathcal C_{01}\) & Yes & Yes & No  & \(240\)\\
\(\mathcal C_{10}\) & Yes & No  & No  & \(240\)\\
\(\mathcal C_{11}\) & Yes & No  & Yes & \(720\)\\
\bottomrule
\end{tabular}
\caption{The four \(S_5\)-invariant binary cover classes over
\(L(P)\).}
\label{tab:invariant-cover-square}
\end{table}

The classes form the affine square

\[
\mathcal C_{00},
\qquad
\mathcal C_{01},
\qquad
\mathcal C_{10},
\qquad
\mathcal C_{11},
\]

with addition inherited from the vector space

\[
H^1(Y,\Ftwo)^{S_5}.
\]

In the earlier descriptive terminology, the two intermediate nonzero
classes were called the native and alternative classes. The coordinate
notation is used here because it records the invariant content directly.

\subsection{The all-one class}

Let

\[
\mathbf 1:E(Y)\longrightarrow\Ftwo
\]

be the constant-one voltage assignment. Its derived graph joins

\[
(u,\varepsilon)
\]

to

\[
(v,\varepsilon+1)
\]

over every edge \(uv\) of \(Y\). Therefore its derived graph is precisely

\[
\KC(Y)=X.
\]

Every triangle and every pentagon has odd length, so

\[
t([\mathbf 1])=1,
\qquad
p([\mathbf 1])=1.
\]

Thus

\[
[\mathbf 1]=\mathcal C_{11}.
\]

\begin{proposition}
\label{prop:x-all-one-class}

The graph \(X\) is the derived graph of the invariant class

\[
\mathcal C_{11}.
\]

Its automorphism order is \(720\), in agreement with
Theorem~\ref{thm:full-automorphism-group}.
\end{proposition}

\subsection{Four classes and three frames}

The four cohomology classes classify invariant covers over a fixed labeled
base graph. The three quotient frames of Section~7 are instead three
covering presentations inside the single total graph \(X\).

% STATUS: Certified alignment across the three quotient presentations
% SOURCE: Project 42 explicit quotient-frame certificate 004

After identifying each quotient with \(L(P)\) and allowing switching of
fiber labels, the three projections

\[
X\longrightarrow X/\langle\tau_i\rangle\cong L(P)
\]

represent the all-one invariant class

\[
\mathcal C_{11}.
\]

Thus the numbers four and three refer to different structures:

\[
\boxed{
\text{four invariant switching classes over one fixed base}
}
\]

and

\[
\boxed{
\text{three conjugate quotient frames inside one total graph}.
}
\]

No claim is made here that the intermediate invariant classes possess only
one quotient presentation, or that multiple quotient frames characterize
\(\mathcal C_{11}\) among all possible total graphs.
